\section{Introduction}\label{sec:intro}

Crash frequency models relate the number of crashes on a roadway entity over a period of time to traffic exposure, geometric design and environmental conditions. They underpin safety performance functions, network screening and the evaluation of countermeasures \citep{lord2010statistical}. Over the past fifteen years, models with random parameters have become the standard way to accommodate unobserved heterogeneity, the fact that the effect of a measured factor varies across sites because of conditions that are not measured \citep{anastasopoulos2009note, mannering2014analytic, mannering2016unobserved}. In a random parameters negative binomial (RPNB) model the coefficient of a covariate is allowed to vary across sites according to a distribution, usually normal, whose mean and standard deviation are estimated from the data, either by simulated maximum likelihood \citep{train2009discrete, bhat2003simulation} or by Markov chain Monte Carlo \citep{elbasyouny2009accident, buddhavarapu2016modeling}.

The flexibility comes with a decision that receives far less attention than the estimation itself: the analyst must choose which coefficients are random. With $p$ candidate covariates, each of which may be excluded, enter with a fixed coefficient, or enter with a random coefficient, there are $3^p$ candidate specifications before heterogeneity in means, heterogeneity in variances or correlation among random parameters is considered. In practice the decision is made sequentially. Covariates are screened in a fixed parameters model, each surviving coefficient is then tried as random, and a random parameter is kept when its estimated standard deviation is statistically significant and fixed otherwise \citep[e.g.,][]{zeng2017multivariate}. More recent work searches the specification space with metaheuristics that minimize an information criterion \citep{ahern2024extensive}. Both strategies settle on a single specification and then report standard errors, marginal effects and policy implications conditional on it, as though it had been fixed before the data were seen.

This practice has three consequences. First, specification uncertainty is not reported, even though it can be large when heterogeneity is modest relative to the Poisson and gamma noise in crash counts. Second, because many standard deviations are tested, often in an order chosen by the analyst, the probability of declaring at least one spurious random parameter grows with the number of candidates and is not controlled. Third, as we show in Section~\ref{sec:identification}, a normally distributed random parameter changes the distribution of counts through a covariate-dependent variance of a lognormal error. That variance adds a quadratic term to the conditional mean and makes overdispersion depend on the covariate, so a detected random parameter may reflect an omitted nonlinear effect or a covariate-dependent dispersion as much as genuine heterogeneity in the effect. Panel data contain information that separates these explanations, but our simulations show that a misspecified model does not use it.

This paper treats the specification of a random parameters model as a variable selection problem and solves it inside a single Bayesian model. Each candidate covariate is assigned one of three states, excluded, fixed or random, and the posterior distribution over these states is computed jointly with the model parameters. The contributions are as follows.
\begin{enumerate}[label=(\roman*),leftmargin=*]
\item \emph{A three-state spike-and-slab prior for effect structure.} Writing the random parameter in non-centered form, $\beta_{ij}=\beta_j+\sigma_j z_{ij}$ with $z_{ij}\sim\Normal(0,1)$ and $\sigma_j$ allowed to take either sign, turns the question ``is this parameter random?'' into the question ``is $\sigma_j$ zero?'', which has the same form as the question ``is $\beta_j$ zero?'' \citep{fruhwirth2010stochastic}. Point-mass spikes on $\beta_j$ and $\sigma_j$, linked by a heredity prior and given beta-binomial hyperpriors, produce posterior probabilities for the three states with an automatic adjustment for the number of candidates \citep{scott2010bayes}.
\item \emph{Identification results that matter for practice.} We characterize how random parameters enter the distribution of crash counts in cross-sectional and panel data, show that the random parameter variance is confounded with the fixed effect for binary covariates and with curvature for continuous covariates, and show that within-site covariance in panel data identifies random parameters separately from the negative binomial dispersion. We then show by simulation that this information is not enough when the mean or the dispersion is misspecified, which motivates selecting the mean curvature and the dispersion covariates jointly with the random parameters.
\item \emph{An efficient sampler.} P\'olya-Gamma augmentation \citep{polson2013bayesian} and the Chinese restaurant table augmentation for the negative binomial dispersion \citep{zhou2015negative} make all regression updates conditionally Gaussian. The step that decides whether a parameter is random integrates out the site-specific draws $z_{ij}$, which reduces it to a one-dimensional integral with a closed-form integrand. Without this step the sampler rarely switches a random parameter on, because a freshly drawn $z_{ij}$ carries no information about the data.
\item \emph{Evidence.} In a simulation study with 100 replications per design, the method recovers the full specification in 60\% of replications with 1{,}000 sites observed over five periods, against 1\% for the sequential procedure, none for its Bayesian counterpart based on the deviance information criterion and 37\% for an information criterion search, and its intervals for the random-parameter standard deviations cover the truth far more often. A false discovery rate rule trades detection for control of false claims of heterogeneity. Omitted curvature and covariate-dependent dispersion produce spurious random parameters under every method that holds the mean linear and the dispersion constant. Squared covariates as fixed-only candidates remove the first. Spike-and-slab selection of dispersion covariates, updated by reversible jump, removes the second in panel data (spurious random parameters in 1\% of replications against 90\% or more) and reduces it from over 90\% to 35\% in cross-sectional data, at little cost when the dispersion is in fact constant. \tbd{Application (Section~\ref{sec:application}).}
\end{enumerate}

The remainder of the paper is organized as follows. Section~\ref{sec:background} reviews random parameters crash models and Bayesian variable selection. Section~\ref{sec:model} presents the model, the identification results and the prior. Section~\ref{sec:computation} develops the sampler. Section~\ref{sec:simulation} reports a simulation study and Section~\ref{sec:application} an application to \tbd{data set}. Section~\ref{sec:discussion} discusses extensions and limitations and Section~\ref{sec:conclusion} concludes. Proofs and derivations are collected in the appendices.
